\documentclass[11pt]{article}
\usepackage{amsmath,amssymb,amsthm,booktabs}
\usepackage[margin=1in]{geometry}

\newcommand{\Psym}{\mathsf P_\infty}
\newcommand{\Cell}{C_\ell}
\newcommand{\Peta}{P_\eta}
\newcommand{\Keta}{\mathcal K_\eta}
\newcommand{\Aeta}{\mathcal A_\eta}

\title{Step 179: Branch A Calkin-Class Attack}
\author{Six Birds Foundations III}
\date{2026-05-15}

\begin{document}
\maketitle

\section{Verdict}

\[
\boxed{\texttt{calkin\_class\_verdict = V\_branch\_a\_stuck\_at\_kappa}.}
\]

Step 173 supplies an operational formula for \(\Psym\), hence for
\[
\Cell=(I-\Psym)M_{m_\ell}\Psym.
\]
But the Branch A Calkin object is
\[
\Cell\Peta,
\]
and the projection \(\Peta\) is onto the actual pulled-evaluator span
\[
H_\eta=\overline{\operatorname{span}\{J_a^*Y^a_{\rho,k}\}}.
\]
By Step 153,
\[
\mathcal U_\infty(J_a^*Y^a_{\rho,k})
=
T_a^*\partial_{\overline\rho}^{\,k}K_a^\Gamma(\cdot,\rho)
=:\kappa_{a,\rho,k}.
\]
Thus deciding the Calkin class requires the same \(\kappa\) record isolated in
Steps 177--178.

\section{T1. Inherited Records}

Step 162 declares
\[
\Aeta=C^*(P_\infty,M_{m_\ell},P_\eta,I),
\qquad
\Keta=\Aeta\cap\mathcal K(H),
\]
and the active quotient object
\[
q_\eta(\Cell\Peta)\in\Aeta/\Keta.
\]

Step 154 gives
\[
\Cell=(I-P_\infty)M_{m_\ell}P_\infty.
\]

Step 158 gives the Calkin dichotomy: \(\Cell\Peta\) is compact if it sends
every normalized weak-null sequence in \(H_\eta\) to norm-null, while positive
essential norm follows from a normalized weak-null sequence
\[
u_n\in H_\eta,\qquad u_n\rightharpoonup0,\qquad
\liminf_n\|\Cell u_n\|>0.
\]

Steps 158--160 also warn that hard-support, Hardy, and finite-window Weyl
evidence is public-shadow or moving-window support only unless promoted to the
actual Burnol/Sonine quotient.

Step 173 gives
\[
(\Psym F)(\tau)
=F(\tau)-\int\frac{\sin(\tau-u)}{\pi(\tau-u)}F(u)\,du
-\sum_n\Psi_n(\tau)\langle F,\Psi_n\rangle.
\]
This acts on known Mellin vectors.

\section{T2a. Matrix Elements}

Let
\[
\kappa_i=\mathcal U_\infty\eta_i.
\]
Then
\[
\langle \eta_i,\Cell\eta_j\rangle
=
\langle \kappa_i,(I-\Psym)M_{m_\ell}\Psym\kappa_j\rangle.
\]
Since \(\Psym\) is orthogonal,
\[
\boxed{
\langle \eta_i,\Cell\eta_j\rangle
=
\langle \kappa_i-\Psym\kappa_i,\,
M_{m_\ell}\Psym\kappa_j\rangle.
}
\]
Step 173 evaluates \(\Psym\kappa_i\) only after \(\kappa_i\) is known.  Thus
the matrix-element route is blocked at \(\kappa\).

\section{T2b. Hilbert-Schmidt and Trace Route}

The Hilbert-Schmidt test would require
\[
\|\Cell\Peta\|_{\mathrm{HS}}^2
=
\operatorname{tr}(\Peta \Cell^*\Cell \Peta).
\]
Equivalently, with an orthonormal frame \(\{e_j\}\) for \(H_\eta\),
\[
\|\Cell\Peta\|_{\mathrm{HS}}^2
=
\sum_j\|\Cell e_j\|^2.
\]
But \(\Peta\) itself is the projection onto the closed span of the vectors
\[
\kappa_{a,\rho,k}=T_a^*\partial_{\overline\rho}^{\,k}
K_a^\Gamma(\cdot,\rho).
\]
Without their Gram kernel or an orthonormalization of that span, the trace
cannot be computed.  K-infinity alone gives \(\Cell\), not \(\Peta\).

\section{T2c. Weyl Sequence}

Step 158's hard-support Weyl sequence shows the natural obstruction mechanism
in a shadow model.  To prove
\[
\|\Cell\Peta\|_{\mathrm e}>0,
\]
one needs actual vectors
\[
u_n\in H_\eta
\]
with persistent \(\Cell\)-mass.  PSWF or hard-support vectors are ambient
vectors unless one proves they lie in, or can be approximated by, the
\(\kappa\)-span.  That is again a \(\kappa\) or projected-Sonine-density
record.

\section{T2d. Step 175--176 Data}

The nonzero Step 175--176 values
\[
L_{\rho,k}(G)=\langle M_\zeta G,\Psym y_{\rho,k}\rangle
\]
show that certain projected zero-evaluator pairings are nonzero.  They do not
identify
\[
\langle \eta_i,\Cell\eta_j\rangle
\]
without the convention chain mapping \(y_{\rho,k}\), \(\eta_{\rho,k}\), and
\(\kappa_{a,\rho,k}\) into one Hilbert model.  That convention chain is exactly
the Step 178 \(\kappa\) record.

\section{Calkin-Class Evidence}

\[
\begin{array}{l|l}
\toprule
\text{criterion} & \text{status}\\
\midrule
\Cell\Peta\in\mathcal K & \text{not decided}\\
\|\Cell\Peta\|_{\mathrm e}>0 & \text{not decided}\\
\text{Hilbert-Schmidt / trace test} & \text{blocked at }\Peta\text{ kernel}\\
\text{matrix entries} & \text{blocked at }\kappa_i,\kappa_j\\
\text{Weyl sequence} & \text{blocked at actual }H_\eta\text{ membership}\\
\bottomrule
\end{array}
\]

The structural finding is that Branch A's matrix route, Branch A's actual Weyl
promotion route, and Branch B's SL164.1 finite matrix route all reduce to the
same projected Sonine kernel / \(\kappa\) record.

\section{Nonclaim Boundary}

This step does not prove \(\Cell\Peta\) compact.  It does not prove positive
essential norm.  It does not construct a Weyl sequence in \(H_\eta\), and it
does not promote public-shadow or finite-window evidence.  It does not close
\(\Xi_{\mathrm{BC}}\) or alter retained no-gos.

\end{document}
